\section{Evaluation}
\label{sec:eval}

We evaluate \sys along five questions:

\begin{itemize}[leftmargin=*]
  \item \textbf{RQ1:} What is the runtime overhead of \sys's core mechanisms and
        observability tools compared to existing approaches?
  \item \textbf{RQ2:} Can \sys-based memory policies improve performance over
        framework-managed offloading and default UVM across realistic workloads?
  \item \textbf{RQ3:} Can \sys-based memory policies improve performance over
        multi-tenant workloads and improve oversubscribe and context switch performance?
  \item \textbf{RQ4:} Can \sys-based scheduling policies provide tangible gains
        in tail latency and throughput in multi-tenant workloads?
  \item \textbf{RQ5:} Can \sys-device side and optimization improve the device policy performance?
\end{itemize}

\subsection{Methodology and Setup}

We evaluated \sys on two machines. Server~A is a single-socket Intel Core
Ultra~9 285K processor (24 cores) with 128~GB DDR5 RAM and an NVIDIA RTX~5090
GPU. Server~B is a single-socket Intel Xeon~6787P (86~cores, 2.0~GHz, 336~MB
LLC) with 768~GB DDR5 RAM (512~GB CXL module) and an NVIDIA H100 GPU. Unless
otherwise noted, we report results on Server~A and average each metric over
10 trials; when averaging speedups we use geometric means. We compare against
NVIDIA's CUPTI and NVTX-based profiler stack, NVBit gpumemtrace, default UVM
memory management, and the native GPU scheduler without \sys-based policies, as
well as framework-managed offloading where applicable.

\subsubsection{Application workloads}
\label{sec:eval-app-workloads}

We use four unmodified GPU applications from pre-release to evaluate \sys under realistic
end-to-end workloads. In all cases, \sys attaches via runtime and kernel driver
instrumentation; application and framework binaries are unchanged and does not require recompilation.

\paragraph{PyTorch.}
We use PyTorch for both inference and training workloads. For RQ1, we run a
ResNet model on ImageNet-sized batches to measure end-to-end inference latency
under \sys and the PyTorch profiler (CUPTI+NVTX). For RQ2, we use PyTorch-based
GNN training under 1.5$\times$ GPU memory oversubscription to study the impact
of programmable UVM policies on epoch time. \sys is loaded as a runtime
extension and instruments CUDA kernels without changes to the PyTorch code,
models, or training loops.

\paragraph{\texttt{llama.cpp}.}
We use the GPT-OSS-120B configuration in \texttt{llama.cpp} to evaluate expert
offload policies (RQ2). The model is approximately 60~GB on an RTX~5090 with
32~GB of device memory, so inference runs in oversubscribed mode with UVM and
the framework's built-in offload support. \sys attaches
transparently at the driver level and
via device-side handlers that tag expert regions as ``hot'' or ``cold'' and
drive prefetch/eviction decisions.

\paragraph{vLLM.}
We evaluate vLLM on a 30B FP8 MoE model to study KV-cache placement and
offloading policies (RQ2) under realistic serving loads. We use a ShareGPT
benchmark with 100 requests and around 30 concurrent requests to stress KV-cache
growth and reuse. \sys runs alongside vLLM with UVM enabled and implements
sequential prefetch and eviction policies for KV-cache segments via callbacks
and GPU handlers.

\paragraph{Faiss.}
We integrate \sys with the GPU backend of Faiss to evaluate vector search with
oversubscribed indices (RQ2). We consider both index build and query workloads
where the index does not fit in device memory and UVM migrates index pages over
PCIe. \sys observes and steers index-segment placement by instrumenting GPU
memory operations and monitoring PCIe traffic, keeping frequently accessed
centroids and posting lists resident while demoting cold segments.

\subsection{RQ1: Runtime Overhead and Observability}

We evaluate the efficiency of \sys's core mechanisms and its observability
capabilities. We first quantify the base runtime overhead and shared state costs,
then demonstrate that \sys enables practical observability tools with
significantly lower overhead than existing frameworks.

\subsubsection{Runtime and Shared State Overhead}

We first quantify the overhead of the \sys runtime itself, independent of any
specific policy logic. We use two microbenchmarks:

\begin{itemize}[leftmargin=*]
  \item a load/store latency benchmark that issues controlled memory accesses
        from GPU kernels, instrumented with \sys and with NVBit gpumemtrace; and
  \item a ResNet inference workload instrumented with \sys-based tracing and
        with the PyTorch profiler (CUPTI+NVTX).
\end{itemize}

\Cref{fig:runtime-overhead} shows per-access latency slowdown under
\sys and NVBit as access sizes increase. \sys adds substantially less latency
than NVBit, remaining low and stable below 128~KB and increasing modestly
thereafter. \Cref{fig:resnet-latency} compares end-to-end ResNet inference
latency distributions: p10 and mean latencies under \sys are slightly better
than with the PyTorch profiler, and p99 latency is comparable. In both cases,
lightweight launch-time injection and warp-level execution of \sys handlers
keep overhead low.

\begin{figure}
\centering
\fbox{\parbox{0.5\textwidth}{\centering [Figure: Runtime Overhead]}}
\caption{Runtime overhead of \sys vs NVBit across varying access sizes.}
\label{fig:runtime-overhead}
\end{figure}

\begin{figure}
\centering
\fbox{\parbox{0.5\textwidth}{\centering [Figure: Resnet Inference Latency Overhead]}}
\caption{ResNet inference latency with \sys tracing vs PyTorch profiler (CUPTI+NVTX).}
\label{fig:resnet-latency}
\end{figure}

Next, we measure the cost of maintaining shared state between CPU and GPU via
\sys maps. We benchmark shared-memory bandwidth and one-way latency on both
servers using a 1~GiB region and GPU-to-CPU and CPU-to-GPU access patterns.
Servers~A and~B deliver 12--18~GB/s over PCIe with 8--15~µs cross-domain
latencies. These costs are small relative to the time scales of the policies we
implement, and our SIMT-aware aggregation and sampling strategies further reduce
the frequency of cross-domain synchronization. We see this in the use cases
below, where policy and observability overhead remains in the low single-digit
percent range.

\subsubsection{Low-overhead GPU Observability}

We evaluate the suite of bcc-style observability tools built on \sys, measuring
accuracy, overhead, and practical utility across diverse workloads.

\paragraph{Thread divergence and workload balance.}
\texttt{kernelretsnoop} attaches to kernel return paths and records per-thread
timestamps. Applied to a matrix multiply kernel with boundary checks, it shows
that threads handling edge cases (5\% of total) incur 600--850~ns longer
execution than interior threads, with timestamps within 10~ns of hardware
performance counters and 2.1\% overhead. Refactoring boundary handling based on
this insight improves overall kernel performance by 8\%. \texttt{threadhist}
reveals load imbalance in a particle simulation with 1M particles across 256
blocks: it identifies a bimodal distribution of iterations per thread;
rebalancing particles based on this information yields an 18\% throughput
improvement with 1.8\% tool overhead.

\paragraph{Launch latency and memory behavior.}
\texttt{launchlate} combines host and device hooks to measure kernel launch
latency. On a pipeline of 50 small kernels (each around 100~µs execution time),
default synchronous execution yields 150--550~µs launch latency per kernel, for
a total runtime of 32~ms. Switching to CUDA graphs, guided by
\texttt{launchlate}, reduces this to 7.2~ms, with less than 0.5\% overhead.
\texttt{mem\_trace} instruments sampled memory operations and recovers per-warp
bank conflict rates within 2\% of hardware counters, with under 3\% end-to-end
overhead across GEMM, SpMV and stencil kernels (\Cref{fig:obs-overhead}).

\begin{figure}
\centering
\fbox{\parbox{\columnwidth}{\centering [Figure: GPU Observability Tools Overhead Comparison]}}
\caption{Overhead of \sys-based observability tools across workloads.}
\label{fig:obs-overhead}
\end{figure}

\paragraph{Comparison to CUPTI and NVBit.}
\Cref{fig:obs-comparison} compares \sys-based tools against CUPTI and
NVBit for common profiling tasks. CUPTI-based profiling incurs 15--25\%
overhead and typically requires kernel restarts; NVBit adds 30--45\%
overhead for similar metrics. \sys tools average 2--3\% overhead and support
live instrumentation without interrupting execution.

\begin{figure}
\centering
\fbox{\parbox{\columnwidth}{\centering [Figure: Observability Tool Comparison]}}
\caption{Comparison of \sys observability tools vs CUPTI and NVBit.}
\label{fig:obs-comparison}
\end{figure}

\subsection{RQ2: Programmable Memory Management}

We now evaluate \sys's memory policy interface on the application workloads
introduced in Section~\ref{sec:eval-app-workloads}, focusing on oversubscribed
setups where memory placement and prefetching dominate performance.

\paragraph{LLM inference: expert and KV-cache offload.}
For the GPT-OSS-120B model in \texttt{llama.cpp}, a \sys policy that triggers a
max-prefetch strategy on GPU-side memory events (prefetching additional experts
when entering an expert region and evicting based on expert-level hotness)
improves throughput by up to 5$\times$ compared to default UVM and the
framework's own offloading, under the oversubscribed configuration described in
Section~\ref{sec:eval-app-workloads}. This policy is expressed as eBPF callbacks
in the driver plus device-side handlers that mark expert regions as
``will-use soon'' or ``cold'' without modifying the application.

On the vLLM 30B FP8 MoE model, we compare vLLM's own CPU-offload mode (no UVM,
\texttt{--cpu-offload-gb 8}) against \sys, which runs vLLM with UVM and a
simple sequential prefetch policy for KV-cache segments. Using the ShareGPT
benchmark from Section~\ref{sec:eval-app-workloads}, \sys improves mean and p99
time-to-first-token by 1.7--2$\times$ and decoding throughput by about
1.3$\times$ over vLLM's framework-managed offload. Simply enabling UVM without
\sys policies performs worse than vLLM's own offload; the \sys configuration is
roughly 2--3$\times$ faster than this UVM-only baseline.

\paragraph{Vector search and GNN training.}
For Faiss, a \sys policy that adaptively prefetches index segments based
on observed PCIe traffic and access streams reduces index build time by about
1.5$\times$ compared to default UVM behavior when the index does not fit in
device memory. Query workloads see similar benefits from keeping frequently
accessed centroids and posting lists resident and demoting cold segments.

For GNN workloads, we run PyTorch-based training under 1.5$\times$
oversubscription of GPU memory. With default UVM, epoch time is around 71~s.
With a simple \sys policy that performs sequential prefetching of adjacency and
feature blocks and FIFO-style eviction, average epoch time drops to about 26~s,
a 2--3$\times$ improvement. For comparison, the same training job with a 0.8$\times$
working set that fits entirely in GPU memory completes an epoch in 15~s; \sys
thus brings the oversubscribed configuration close to the non-oversubscribed
baseline. The CPU-side framework overhead for these policies is around 1\%, and
device-side instrumentation adds 2--3\%.

Across these applications, the common pattern is that programmable region-based
placement and prefetch policies in \sys allow the system to exploit workload
structure (experts, KV-cache, graph structure, index layout) that existing UVM
policies and framework-level offload mechanisms cannot leverage without invasive
changes.

\subsection{RQ3: Fine-grain Scheduling}

We finally evaluate \sys's scheduling interface using priority-based policies on
mixed workloads that combine high-priority latency-sensitive kernels with
low-priority batch jobs. We focus on tail latency and fairness; power-aware
control is omitted here.

\paragraph{Priority-based preemption.}
We compare three approaches: (1) no prioritization (FIFO), (2) CUDA stream-level
priorities, and (3) \sys-based priority scheduling with warp-level preemption.
The workload combines real-time video processing kernels as high-priority tasks
with batch ML training as low-priority background load.

\Cref{fig:p99-latency} plots p99 latency for high-priority tasks. FIFO exhibits
25--80~ms p99 latency due to blocking behind long-running kernels. CUDA streams
reduce p99 to 15--30~ms but still suffer head-of-line blocking. \sys cuts p99
latency to 5--12~ms by preempting low-priority warps at synchronization points
and admitting high-priority kernels promptly. Scheduling decisions take
0.3--0.8~µs with 4--6\% overhead.

\begin{figure}
\centering
\fbox{\parbox{\columnwidth}{\centering [Figure: Priority Scheduler Tail Latency]}}
\caption{p99 latency for high-priority tasks under different scheduling policies.}
\label{fig:p99-latency}
\end{figure}

\Cref{fig:lowprio-throughput} shows the impact on low-priority throughput.
\sys incurs an 8--12\% throughput reduction compared to FIFO due to preemption
overheads, but maintains better fairness than CUDA streams, which reduce
low-priority throughput by 18--25\% and can starve low-priority work for long
intervals.

\begin{figure}
\centering
\fbox{\parbox{\columnwidth}{\centering [Figure: Low-Priority Task Throughput]}}
\caption{Throughput impact on low-priority tasks with different scheduling policies.}
\label{fig:lowprio-throughput}
\end{figure}

We also applied a \sys-based block scheduler driven by the GPU-side interface
(Section~\ref{sec:design-sched}) to a minimal LLM inference framework. Using a
cluster-style launch where persistent worker blocks pull work units based on
tenant-aware priorities, we observe around 10\% end-to-end improvement in tokens
per second compared to a static block assignment, with similar overheads. Early
experiments on graph algorithms show potential for larger gains when block-level
load imbalance is high, but we leave a full exploration to future work.

Overall, scheduling policies in \sys provide clear tail-latency and fairness
benefits in mixed workloads, while the largest performance gains in our study
come from programmable memory management.
